% 【拆分生成】 answers.tex —— 2021 年·新高考Ⅰ卷 · 参考答案与详细解析片段
% 【拆分生成】 由 tools/split_project.py 自合并书 answers/2021_xgk1_sol.tex 包裹 answer 环境生成，勿手改
% 由 tools/md2tex.py 自 2021_数学真题_含解析.md 生成，勿手改（改源 MD 后重跑）


\begin{gknotice}
\textbf{解析归属}：本文档全部解析由本库编写者根据试题独立推导，\textbf{非官方标准答案}。推导过程完整给出以便核验；若与官方发布版本不一致，以官方为准。

\textbf{历史意义}：这是湖南首次使用新高考 I 卷的年份。
\end{gknotice}

\par\medskip

\gkpart{一、单选题（每小题 5 分，共 40 分）}

\begin{answer}
\gksolhead{1}{}

\textbf{题目}　设集合 $A=\{x\mid-2<x<4\}$，$B=\{2,3,4,5\}$，则 $A\cap B=$（　　）

\textbf{解析}

$A$ 是\textbf{开区间} $(-2,4)$，逐个检验 $B$ 中元素是否满足 $-2<x<4$：

\begin{center}
\begin{tabular}{ccccc}
\toprule
元素 & $2$ & $3$ & $4$ & $5$ \\
\midrule
是否满足 &  &  & （$4$ 不小于 $4$） &  \\
\bottomrule
\end{tabular}
\end{center}

\[A\cap B=\{2,3\}\]

\gkanswer{答案：B}

\begin{gktip}
\textbf{易错点}：$4$ 是区间的\textbf{开端点}，不属于 $A$。选项 D 正是为忽略"开区间不含端点"的学生设置的。
\end{gktip}

\par\medskip

\end{answer}
\begin{answer}
\gksolhead{2}{}

\textbf{题目}　已知 $z=2-\mathrm i$，则 $z(\bar z+\mathrm i)=$（　　）

\textbf{解析}

\textbf{第一步，求共轭。}
\[\bar z=2+\mathrm i\]

\textbf{第二步，算括号。}
\[\bar z+\mathrm i=2+\mathrm i+\mathrm i=2+2\mathrm i\]

\textbf{第三步，乘法展开}（注意 $\mathrm i^{2}=-1$）：
\[z(\bar z+\mathrm i)=(2-\mathrm i)(2+2\mathrm i)=4+4\mathrm i-2\mathrm i-2\mathrm i^{2}=4+2\mathrm i+2=6+2\mathrm i\]

\gkanswer{答案：C}

\par\medskip

\end{answer}
\begin{answer}
\gksolhead{3}{}

\textbf{题目}　已知圆锥的底面半径为 $\sqrt2$，其侧面展开图为一个半圆，则该圆锥的母线长为（　　）

\textbf{解析}

\textbf{核心等式}：圆锥侧面展开后是一个扇形，\textbf{扇形的弧长等于底面圆的周长}，\textbf{扇形的半径等于圆锥的母线}。

设母线长为 $l$。

\textbf{底面周长}：
\[C=2\pi r=2\pi\sqrt2\]

\textbf{展开图（半圆）的弧长}：半圆的半径是 $l$，弧长为整圆周长的一半：
\[\frac12\times2\pi l=\pi l\]

\textbf{令二者相等}：
\[\pi l=2\pi\sqrt2\ \Longrightarrow\ l=2\sqrt2\]

\gkanswer{答案：B}

\begin{gktip}
\textbf{一般结论}：若侧面展开图的圆心角为 $\theta$（弧度），则 $\theta\cdot l=2\pi r$，即 $\dfrac rl=\dfrac\theta{2\pi}$。本题 $\theta=\pi$（半圆），故 $\dfrac rl=\dfrac12$，母线是半径的 $2$ 倍。记住这个比例关系可以口算。
\end{gktip}

\par\medskip

\end{answer}
\begin{answer}
\gksolhead{4}{}

\textbf{题目}　下列区间中，函数 $f(x)=7\sin\left(x-\dfrac\pi6\right)$ 单调递增的区间是（　　）

\textbf{解析}

\textbf{第一步，写出递增区间的一般式。}

$y=\sin u$ 在 $u\in\left[-\dfrac\pi2+2k\pi,\ \dfrac\pi2+2k\pi\right]$ 上递增。令 $u=x-\dfrac\pi6$：
\[-\frac\pi2+2k\pi\le x-\frac\pi6\le\frac\pi2+2k\pi\]

\[-\frac\pi3+2k\pi\le x\le\frac{2\pi}3+2k\pi\]

\textbf{第二步，写出前两个具体区间（便于比对选项）。}

\begin{itemize}
\item $k=0$：$\left[-\dfrac\pi3,\ \dfrac{2\pi}3\right]$，即约 $[-1.05,\ 2.09]$
\item $k=1$：$\left[\dfrac{5\pi}3,\ \dfrac{8\pi}3\right]$，即约 $[5.24,\ 8.38]$
\end{itemize}

\textbf{第三步，逐项检验（选项必须完全落在某个递增区间内）。}

\begin{center}
\begin{tabular}{clll}
\toprule
选项 & 区间 & 近似 & 是否含于递增区间 \\
\midrule
A & $\left(0,\dfrac\pi2\right)$ & $(0,\ 1.57)$ &  含于 $[-1.05,2.09]$ \\
B & $\left(\dfrac\pi2,\pi\right)$ & $(1.57,\ 3.14)$ &  超出 $2.09$ \\
C & $\left(\pi,\dfrac{3\pi}2\right)$ & $(3.14,\ 4.71)$ &  落在递减段 \\
D & $\left(\dfrac{3\pi}2,2\pi\right)$ & $(4.71,\ 6.28)$ &  起点 $4.71<5.24$ \\
\bottomrule
\end{tabular}
\end{center}

\gkanswer{答案：A}

\begin{gktip}
\textbf{方法总结}：这类题不必逐个选项代入验证单调性，\textbf{先把递增区间的具体范围算出来（用小数近似），再看哪个选项是它的子集}，最快。注意选项 D 只是"部分落在"递增区间内，不满足"在该区间上单调递增"的要求。
\end{gktip}

\par\medskip

\end{answer}
\begin{answer}
\gksolhead{5}{}

\textbf{题目}　已知 $F_1$，$F_2$ 是椭圆 $C:\dfrac{x^{2}}9+\dfrac{y^{2}}4=1$ 的两个焦点，点 $M$ 在 $C$ 上，则 $\left|MF_1\right|\cdot\left|MF_2\right|$ 的最大值为（　　）

\textbf{解析}

\textbf{第一步，抓住椭圆定义提供的"和为定值"。}

由 $a^{2}=9$ 得 $a=3$，椭圆定义给出
\[\left|MF_1\right|+\left|MF_2\right|=2a=6\]

\textbf{第二步，"和定则积有最大值"——基本不等式。}

对任意正数 $m$、$n$：
\[mn\le\left(\frac{m+n}2\right)^{2}\]

代入：
\[\left|MF_1\right|\cdot\left|MF_2\right|\le\left(\frac62\right)^{2}=9\]

\textbf{第三步，检验等号能否取到。}

等号成立当且仅当 $\left|MF_1\right|=\left|MF_2\right|=3$。

这意味着 $M$ 到两焦点等距，即 $M$ 在 $y$ 轴上——正是椭圆的\textbf{短轴端点} $(0,\pm2)$。

验证：$b^{2}=4$，$c^{2}=9-4=5$，短轴端点到焦点的距离
\[\sqrt{b^{2}+c^{2}}=\sqrt{4+5}=3\quad\checkmark\]

等号可以取到。

\gkanswer{答案：C}

\begin{gktip}
\textbf{方法总结}：椭圆中出现"$\left|MF_1\right|\cdot\left|MF_2\right|$"、"$\left|MF_1\right|^{2}+\left|MF_2\right|^{2}$"这类对称式，第一反应就是用定义把\textbf{和}换成 $2a$，再用基本不等式或完全平方公式处理。这是圆锥曲线中最常用的一类转化。
\end{gktip}

\par\medskip

\end{answer}
\begin{answer}
\gksolhead{6}{}

\textbf{题目}　若 $\tan\theta=-2$，则 $\dfrac{\sin\theta(1+\sin2\theta)}{\sin\theta+\cos\theta}=$（　　）

\textbf{解析}

\textbf{第一步（本题的题眼）：把 $1+\sin2\theta$ 配成完全平方。}
\[1+\sin2\theta=\underbrace{\sin^{2}\theta+\cos^{2}\theta}_{=1}+2\sin\theta\cos\theta=(\sin\theta+\cos\theta)^{2}\]

\textbf{第二步，约分。}
\[\frac{\sin\theta\cdot(\sin\theta+\cos\theta)^{2}}{\sin\theta+\cos\theta}=\sin\theta(\sin\theta+\cos\theta)=\sin^{2}\theta+\sin\theta\cos\theta\]

\textbf{第三步，齐次化}（分母补上 $\sin^{2}\theta+\cos^{2}\theta=1$，使分子分母同为二次齐次式）：
\[\sin^{2}\theta+\sin\theta\cos\theta=\frac{\sin^{2}\theta+\sin\theta\cos\theta}{\sin^{2}\theta+\cos^{2}\theta}\]

分子分母同除以 $\cos^{2}\theta$：
\[=\frac{\tan^{2}\theta+\tan\theta}{\tan^{2}\theta+1}\]

\textbf{第四步，代入 $\tan\theta=-2$。}
\[=\frac{(-2)^{2}+(-2)}{(-2)^{2}+1}=\frac{4-2}{4+1}=\frac25\]

\gkanswer{答案：C}

\begin{gktip}
\textbf{方法总结}：三角求值题的两个高频技巧本题都用到了——
① \textbf{$1\pm\sin2\theta=(\sin\theta\pm\cos\theta)^{2}$}（"1"的代换 + 完全平方），用于约分；
② \textbf{齐次化}：只知道 $\tan\theta$ 时，把关于 $\sin$、$\cos$ 的\textbf{齐次式}（分子分母次数相同）同除以 $\cos^{n}\theta$，全部转成 $\tan\theta$。注意齐次化前必须先保证分子分母次数相同，本题通过给分子补分母 $1$ 实现。
\end{gktip}

\par\medskip

\end{answer}
\begin{answer}
\gksolhead{7}{}

\textbf{题目}　若过点 $(a,b)$ 可以作曲线 $y=\mathrm e^{x}$ 的两条切线，则（　　）

\textbf{解析}

\textbf{第一步，写出"过定点的切线"条件。}

设切点为 $\left(t,\mathrm e^{t}\right)$。由 $y'=\mathrm e^{x}$ 知切线斜率为 $\mathrm e^{t}$，切线方程：
\[y-\mathrm e^{t}=\mathrm e^{t}(x-t)\]

过点 $(a,b)$：
\[b-\mathrm e^{t}=\mathrm e^{t}(a-t)\ \Longrightarrow\ b=\mathrm e^{t}(a-t)+\mathrm e^{t}=\mathrm e^{t}(a+1-t)\]

\textbf{"有两条切线"等价于方程 $b=h(t)$ 有两个不同的实根}，其中
\[h(t)=\mathrm e^{t}(a+1-t)\]

\textbf{第二步，研究 $h(t)$ 的图象。}
\[h'(t)=\mathrm e^{t}(a+1-t)+\mathrm e^{t}\cdot(-1)=\mathrm e^{t}(a-t)\]

因 $\mathrm e^{t}>0$，符号由 $a-t$ 决定：

\begin{center}
\begin{tabular}{lccc}
\toprule
$t$ & $(-\infty,a)$ & $a$ & $(a,+\infty)$ \\
\midrule
$h'$ & $+$ & $0$ & $-$ \\
$h$ & 增 & 极大 & 减 \\
\bottomrule
\end{tabular}
\end{center}

\textbf{最大值}：
\[h(a)=\mathrm e^{a}(a+1-a)=\mathrm e^{a}\]

\textbf{两端的极限}：
\begin{itemize}
\item $t\to-\infty$：$\mathrm e^{t}\to0$ 而 $(a+1-t)\to+\infty$，但指数衰减占优，故 $h(t)\to0^{+}$
\item $t\to+\infty$：$\mathrm e^{t}\to+\infty$ 而 $(a+1-t)\to-\infty$，故 $h(t)\to-\infty$
\end{itemize}

\textbf{第三步，数交点。}

\begin{itemize}
\item \textbf{左分支}（$t<a$）：从 $0^{+}$ 单调增到 $\mathrm e^{a}$，值域 $\left(0,\mathrm e^{a}\right)$
\item \textbf{右分支}（$t>a$）：从 $\mathrm e^{a}$ 单调减到 $-\infty$，值域 $\left(-\infty,\mathrm e^{a}\right)$
\end{itemize}

水平线 $y=b$ 要与图象有\textbf{两个}交点，必须同时被两个分支取到，即 $b$ 落在两值域的交集：
\[\left(0,\mathrm e^{a}\right)\cap\left(-\infty,\mathrm e^{a}\right)=\left(0,\mathrm e^{a}\right)\]

\[0<b<\mathrm e^{a}\]

\gkanswer{答案：D}

\begin{gktip}
\textbf{几何直观}：$y=\mathrm e^{x}$ 的图象整体在 $x$ 轴上方且是凸函数。要能作出\textbf{两条}切线，点 $(a,b)$ 必须位于"曲线下方、$x$ 轴上方"的区域。$b\le0$ 时只能作一条（或零条），$b\ge\mathrm e^{a}$ 时点在曲线上或上方，也作不出两条。选项 D 正是这个区域的代数描述。
\end{gktip}

\par\medskip

\end{answer}
\begin{answer}
\gksolhead{8}{}

\textbf{题目}　$6$ 个球标有 $1\sim6$，有放回随机取两次。甲："第一次是 $1$"；乙："第二次是 $2$"；丙："两次之和是 $8$"；丁："两次之和是 $7$"。则（　　）

\textbf{解析}

\textbf{样本空间}：有放回取两次，共 $6\times6=36$ 个等可能结果。

\textbf{先算各事件的概率}：
\[P(\text{甲})=\frac6{36}=\frac16,\qquad P(\text{乙})=\frac6{36}=\frac16\]

和为 $8$ 的结果：$(2,6),(3,5),(4,4),(5,3),(6,2)$ 共 $5$ 个：
\[P(\text{丙})=\frac5{36}\]

和为 $7$ 的结果：$(1,6),(2,5),(3,4),(4,3),(5,2),(6,1)$ 共 $6$ 个：
\[P(\text{丁})=\frac6{36}=\frac16\]

\textbf{逐项用独立性定义 $P(XY)=P(X)P(Y)$ 检验：}

\par\medskip

\textbf{A（甲与丙）}

甲$\cap$丙 = "第一次为 $1$ 且和为 $8$" → 第二次需为 $7$，\textbf{不可能}：
\[P(\text{甲}\cap\text{丙})=0\]

而
\[P(\text{甲})P(\text{丙})=\frac16\times\frac5{36}=\frac5{216}\ne0\]

不独立。

\par\medskip

\textbf{B（甲与丁）}

甲$\cap$丁 = "第一次为 $1$ 且和为 $7$" → 第二次为 $6$，恰好 $1$ 个结果：
\[P(\text{甲}\cap\text{丁})=\frac1{36}\]

而
\[P(\text{甲})P(\text{丁})=\frac16\times\frac16=\frac1{36}\quad\checkmark\]

\textbf{相互独立} 

\par\medskip

\textbf{C（乙与丙）}

乙$\cap$丙 = "第二次为 $2$ 且和为 $8$" → 第一次为 $6$，$1$ 个结果：
\[P(\text{乙}\cap\text{丙})=\frac1{36}=\frac6{216}\]

而
\[P(\text{乙})P(\text{丙})=\frac16\times\frac5{36}=\frac5{216}\]

$\dfrac6{216}\ne\dfrac5{216}$，不独立。

\par\medskip

\textbf{D（丙与丁）}

和不可能既是 $8$ 又是 $7$，两事件\textbf{互斥}：
\[P(\text{丙}\cap\text{丁})=0\ne P(\text{丙})P(\text{丁})\]

不独立。

\gkanswer{答案：B}

\begin{gktip}
\textbf{关键洞察}：为什么只有"和为 $7$"与单次结果独立？因为 $7=1+6=2+5=\cdots$，\textbf{每个首次取值都恰好有一个第二次取值使和为 $7$}，所以知道第一次的结果不改变"和为 $7$"的概率（始终是 $\dfrac16$）。而和为 $8$ 时，首次取 $1$ 就绝无可能，首次取 $4$ 却有 $\dfrac16$ 的机会——信息发生了改变，故不独立。

\textbf{另一个易错点}：\textbf{互斥 $\neq$ 独立}，事实上两个概率非零的事件若互斥，则一定\textbf{不}独立（D 项）。
\end{gktip}

\par\medskip

\end{answer}
\gkpart{二、多选题（每小题 5 分，共 20 分）}

\begin{answer}
\gksolhead{9}{}

\textbf{题目}　样本数据 $x_1,\cdots,x_n$ 与 $y_i=x_i+c$（$c$ 为非零常数），则（　　）

\textbf{解析}

这是\textbf{整体平移}变换，把每个数据都加上同一个常数 $c$。

\textbf{A（平均数）}
\[\bar y=\frac1n\sum\left(x_i+c\right)=\bar x+c\ne\bar x\quad(c\ne0)\]

\textbf{B（中位数）}

平移不改变数据的大小顺序，故中位数也整体平移 $c$，不相同。

\textbf{C（标准差）}

标准差衡量数据\textbf{相对于自身平均数的偏离程度}。因 $y_i-\bar y=\left(x_i+c\right)-\left(\bar x+c\right)=x_i-\bar x$，每个偏差都不变，故
\[s_y=\sqrt{\frac1n\sum\left(y_i-\bar y\right)^{2}}=\sqrt{\frac1n\sum\left(x_i-\bar x\right)^{2}}=s_x\quad\checkmark\]

\textbf{D（极差）}
\[y_{\max}-y_{\min}=\left(x_{\max}+c\right)-\left(x_{\min}+c\right)=x_{\max}-x_{\min}\quad\checkmark\]

\gkanswer{答案：CD}

\begin{gktip}
\textbf{通用规律}（值得背下来）：对线性变换 $y=kx+c$，
\begin{itemize}
\item \textbf{位置类}统计量（平均数、中位数、众数、百分位数）：$\to k\cdot(\text{\text{原值}})+c$
\item \textbf{离散类}统计量（标准差、极差）：$\to|k|\cdot(\text{\text{原值}})$，\textbf{与 $c$ 无关}
\item \textbf{方差}：$\to k^{2}\cdot(\text{\text{原方差}})$
\end{itemize}

本题 $k=1$，故离散类全部不变、位置类全部平移。
\end{gktip}

\par\medskip

\end{answer}
\begin{answer}
\gksolhead{10}{}

\textbf{题目}　$O$ 为原点，$P_1(\cos\alpha,\sin\alpha)$，$P_2(\cos\beta,-\sin\beta)$，$P_3(\cos(\alpha+\beta),\sin(\alpha+\beta))$，$A(1,0)$，则（　　）

\textbf{解析}

\textbf{先观察}：$P_1$、$P_2$、$P_3$、$A$ 都在\textbf{单位圆}上（每个点的坐标平方和为 $1$）。

\par\medskip

\textbf{A 项 }
\[\left|\overrightarrow{OP_1}\right|=\sqrt{\cos^{2}\alpha+\sin^{2}\alpha}=1,\qquad \left|\overrightarrow{OP_2}\right|=\sqrt{\cos^{2}\beta+\sin^{2}\beta}=1\]

两者都是单位圆半径，相等 

\par\medskip

\textbf{B 项 }
\[\left|\overrightarrow{AP_1}\right|^{2}=(\cos\alpha-1)^{2}+\sin^{2}\alpha=\cos^{2}\alpha-2\cos\alpha+1+\sin^{2}\alpha=2-2\cos\alpha\]

\[\left|\overrightarrow{AP_2}\right|^{2}=(\cos\beta-1)^{2}+(-\sin\beta)^{2}=2-2\cos\beta\]

两者相等当且仅当 $\cos\alpha=\cos\beta$，\textbf{不恒成立}（例如 $\alpha=0$、$\beta=\dfrac\pi2$）

\par\medskip

\textbf{C 项 }

\textbf{左边}：
\[\overrightarrow{OA}\cdot\overrightarrow{OP_3}=1\times\cos(\alpha+\beta)+0\times\sin(\alpha+\beta)=\cos(\alpha+\beta)\]

\textbf{右边}：
\[\overrightarrow{OP_1}\cdot\overrightarrow{OP_2}=\cos\alpha\cos\beta+\sin\alpha\cdot(-\sin\beta)=\cos\alpha\cos\beta-\sin\alpha\sin\beta\]

由\textbf{两角和的余弦公式}，右边 $=\cos(\alpha+\beta)$，与左边相等 

\begin{gktip}
这个等式的几何意义：$P_2$ 的纵坐标取了负号（相当于 $-\beta$ 角），因此 $\overrightarrow{OP_1}$ 与 $\overrightarrow{OP_2}$ 的夹角恰为 $\alpha-(-\beta)=\alpha+\beta$，而两个单位向量的数量积就是夹角余弦。
\end{gktip}

\par\medskip

\textbf{D 项 }

\textbf{左边}：
\[\overrightarrow{OA}\cdot\overrightarrow{OP_1}=\cos\alpha\]

\textbf{右边}：
\[\overrightarrow{OP_2}\cdot\overrightarrow{OP_3}=\cos\beta\cos(\alpha+\beta)+(-\sin\beta)\sin(\alpha+\beta)\]
\[=\cos\beta\cos(\alpha+\beta)-\sin\beta\sin(\alpha+\beta)=\cos\left[\beta+(\alpha+\beta)\right]=\cos(\alpha+2\beta)\]

$\cos(\alpha+2\beta)$ 与 $\cos\alpha$ 一般不等（如 $\alpha=0$、$\beta=\dfrac\pi2$ 时左 $=1$、右 $=\cos\pi=-1$）

\gkanswer{答案：AC}

\begin{gktip}
\textbf{方法总结}：单位圆上点的数量积题，核心是"\textbf{两个单位向量的数量积 = 夹角的余弦}"。做题时把每个点对应的\textbf{角}标清楚（注意 $P_2$ 的角是 $-\beta$），大部分选项可以直接由夹角关系判断，不必展开计算。
\end{gktip}

\par\medskip

\end{answer}
\begin{answer}
\gksolhead{11}{}

\textbf{题目}　点 $P$ 在圆 $(x-5)^{2}+(y-5)^{2}=16$ 上，$A(4,0)$，$B(0,2)$，则（　　）

\textbf{解析}

圆心 $C(5,5)$，半径 $r=4$。

\par\medskip

\textbf{A、B 项（点到直线的距离范围）}

\textbf{直线 $AB$ 的方程}（用截距式）：
\[\frac x4+\frac y2=1\ \Longrightarrow\ x+2y-4=0\]

\textbf{圆心到 $AB$ 的距离}：
\[d=\frac{|5+2\times5-4|}{\sqrt{1^{2}+2^{2}}}=\frac{11}{\sqrt5}=\frac{11\sqrt5}5\approx4.92\]

因 $d\approx4.92>r=4$，\textbf{直线与圆相离}，故圆上点到直线的距离范围为
\[[d-r,\ d+r]=[4.92-4,\ 4.92+4]=[0.92,\ 8.92]\]

\begin{itemize}
\item \textbf{A }：最大值 $\approx8.92<10$，故所有点的距离都小于 $10$ 
\item \textbf{B }：最小值 $\approx0.92<2$，存在点的距离不大于 $2$ 
\end{itemize}

\par\medskip

\textbf{C、D 项（视角的最值）}

\textbf{从 $B$ 看圆的切线长}：
\[|BC|=\sqrt{(5-0)^{2}+(5-2)^{2}}=\sqrt{25+9}=\sqrt{34}\]

\[\text{切线长}=\sqrt{|BC|^{2}-r^{2}}=\sqrt{34-16}=\sqrt{18}=3\sqrt2\]

\textbf{关键分析}：当 $P$ 在圆上移动时，射线 $BP$ 的方向在\textbf{从 $B$ 出发的两条切线之间}摆动。

\begin{itemize}
\item 从 $B$ 看圆的\textbf{半张角} $\varphi$ 满足 $\sin\varphi=\dfrac r{|BC|}=\dfrac4{\sqrt{34}}\approx0.686$，故 $\varphi\approx43.3^\circ$
\item 射线 $BA$ 与 $BC$ 的夹角：$\overrightarrow{BA}=(4,-2)$、$\overrightarrow{BC}=(5,3)$，
\end{itemize}
\[\cos\angle ABC=\frac{20-6}{\sqrt{20}\times\sqrt{34}}\approx0.537\ \Longrightarrow\ \angle ABC\approx57.5^\circ\]

因 $\angle ABC\approx57.5^\circ>\varphi\approx43.3^\circ$，说明\textbf{射线 $BA$ 落在圆的视角范围之外}（$BA$ 不穿过圆）。

因此 $\angle PBA$ 的取值范围是
\[\left[\angle ABC-\varphi,\ \angle ABC+\varphi\right]\approx\left[14.2^\circ,\ 100.8^\circ\right]\]

\textbf{两个端点都在切线位置取到}，而此时
\[|PB|=\text{切线长}=3\sqrt2\]

\begin{itemize}
\item \textbf{C }：$\angle PBA$ 最小时 $P$ 是靠近 $BA$ 一侧的切点，$|PB|=3\sqrt2$ 
\item \textbf{D }：$\angle PBA$ 最大时 $P$ 是另一侧的切点，$|PB|=3\sqrt2$ 
\end{itemize}

\gkanswer{答案：ACD}

\begin{gktip}
\textbf{方法总结}：
① "圆上点到定直线的距离范围"永远是 $[d-r,\ d+r]$（相离时），这是标准模型；
② "从圆外一点看圆上动点的视角最值"永远在\textbf{切点}处取到，此时该点到动点的距离就是\textbf{切线长} $\sqrt{|BC|^{2}-r^{2}}$。

C、D 同时正确并不矛盾——最大最小分别对应两个不同的切点，但它们到 $B$ 的距离\textbf{相同}（切线长唯一）。
\end{gktip}

\par\medskip

\end{answer}
\begin{answer}
\gksolhead{12}{}

\textbf{题目}　正三棱柱 $ABC\text{-}A_1B_1C_1$ 中 $AB=AA_1=1$，$\overrightarrow{BP}=\lambda\overrightarrow{BC}+\mu\overrightarrow{BB_1}$，$\lambda,\mu\in[0,1]$，则（　　）

\textbf{解析}

\textbf{建系}：底面是边长 $1$ 的正三角形，高为 $\dfrac{\sqrt3}2$。以 $B$ 为原点，$BC$ 为 $x$ 轴、竖直方向为 $z$ 轴：
\[B(0,0,0),\quad C(1,0,0),\quad A\left(\frac12,\frac{\sqrt3}2,0\right)\]
\[B_1(0,0,1),\quad C_1(1,0,1),\quad A_1\left(\frac12,\frac{\sqrt3}2,1\right)\]

由 $\overrightarrow{BP}=\lambda(1,0,0)+\mu(0,0,1)$：
\[P(\lambda,\ 0,\ \mu)\]

\textbf{注意}：$P$ 的 $y$ 坐标恒为 $0$，说明 $P$ 在\textbf{侧面 $BCC_1B_1$}（一个边长为 $1$ 的正方形）上移动。

\par\medskip

\textbf{A 项 }

$\lambda=1$ 时 $P(1,0,\mu)$，即 $P$ 在棱 $CC_1$ 上移动。

$\triangle AB_1P$ 的周长 $=\left|AB_1\right|+\left|B_1P\right|+|PA|$。其中
\[\left|B_1P\right|=\sqrt{(1-0)^{2}+0+(\mu-1)^{2}}=\sqrt{1+(\mu-1)^{2}}\]

$\mu=0$ 时为 $\sqrt2$，$\mu=1$ 时为 $1$——\textbf{随 $\mu$ 变化}，故周长不是定值 

\par\medskip

\textbf{B 项 }

$\mu=1$ 时 $P(\lambda,0,1)$，即 $P$ 在棱 $B_1C_1$ 上移动。

三棱锥 $P\text{-}A_1BC$ 以 $\triangle A_1BC$ 为底面（\textbf{固定}），体积取决于 $P$ 到平面 $A_1BC$ 的距离。

因 $B_1C_1\parallel BC$（正三棱柱的对应棱平行），而 $BC\subset$ 平面 $A_1BC$，故
\[B_1C_1\parallel\text{平面 }A_1BC\]

\textbf{平行线上任一点到该平面的距离相等}，所以 $P$ 到平面 $A_1BC$ 的距离为定值，体积为定值 

\par\medskip

\textbf{C 项 }

$\lambda=\dfrac12$ 时 $P\left(\dfrac12,0,\mu\right)$。
\[\overrightarrow{A_1P}=P-A_1=\left(0,\ -\frac{\sqrt3}2,\ \mu-1\right),\qquad \overrightarrow{BP}=P-B=\left(\frac12,\ 0,\ \mu\right)\]

由 $A_1P\perp BP$：
\[\overrightarrow{A_1P}\cdot\overrightarrow{BP}=0\times\frac12+\left(-\frac{\sqrt3}2\right)\times0+(\mu-1)\mu=\mu^{2}-\mu=0\]

\[\mu(\mu-1)=0\ \Longrightarrow\ \mu=0\ \text{或}\ \mu=1\]

\textbf{两个解}（对应 $P=B$ 和 $P$ 在 $B_1C_1$ 的中点上方），并非"有且仅有一个" 

\par\medskip

\textbf{D 项 }

$\mu=\dfrac12$ 时 $P\left(\lambda,0,\dfrac12\right)$。
\[\overrightarrow{A_1B}=B-A_1=\left(-\frac12,\ -\frac{\sqrt3}2,\ -1\right)\]

要 $A_1B\perp$ 平面 $AB_1P$，需 $\overrightarrow{A_1B}$ 与平面内\textbf{两条相交向量}都垂直。

\textbf{先验证 $\overrightarrow{AB_1}$}：
\[\overrightarrow{AB_1}=B_1-A=\left(-\frac12,\ -\frac{\sqrt3}2,\ 1\right)\]
\[\overrightarrow{A_1B}\cdot\overrightarrow{AB_1}=\frac14+\frac34-1=0\quad\checkmark\]

（这一条\textbf{自动成立}，与 $\lambda$ 无关。）

\textbf{再看 $\overrightarrow{AP}$}：
\[\overrightarrow{AP}=P-A=\left(\lambda-\frac12,\ -\frac{\sqrt3}2,\ \frac12\right)\]

\[\overrightarrow{A_1B}\cdot\overrightarrow{AP}=-\frac12\left(\lambda-\frac12\right)+\frac34-\frac12=-\frac\lambda2+\frac14+\frac14=\frac12-\frac\lambda2\]

令其为零：
\[\lambda=1\]

\textbf{唯一解}（此时 $P=C_1$ 的中点位置，即 $P\left(1,0,\dfrac12\right)$，在棱 $CC_1$ 的中点）

\gkanswer{答案：BD}

\begin{gktip}
\textbf{方法总结}：立体几何多选题中含参动点时，\textbf{建系后把 $P$ 的坐标用参数表示}是最稳妥的路线。本题四个选项分别考查"周长是否定值""体积是否定值""垂直方程的解的个数"三类问题，都可以在坐标系下机械地判断。

特别注意 C、D 两项问的是"\textbf{有且仅有一个}点"——必须解出方程的\textbf{全部}根再计数，C 项的失分点正是漏掉 $\mu=0$ 这个根。
\end{gktip}

\par\medskip

\end{answer}
\gkpart{三、填空题（每小题 5 分，共 20 分）}

\begin{answer}
\gksolhead{13}{}

\textbf{题目}　已知函数 $f(x)=x^{3}\left(a\cdot2^{x}-2^{-x}\right)$ 是偶函数，则 $a=$ \gkblank{}。

\textbf{解析}

\textbf{方法一（利用奇偶性的乘法规律，最快）}

$f(x)$ 是两个因子的乘积：$x^{3}$ 与 $g(x)=a\cdot2^{x}-2^{-x}$。

已知 $x^{3}$ 是\textbf{奇函数}，而 $f$ 是\textbf{偶函数}。由"奇 × 奇 = 偶"，可知 $g(x)$ 必须是\textbf{奇函数}。

奇函数在 $x=0$ 处有定义时必有 $g(0)=0$：
\[g(0)=a\cdot2^{0}-2^{0}=a-1=0\ \Longrightarrow\ a=1\]

\textbf{验证}：$a=1$ 时 $g(x)=2^{x}-2^{-x}$，$g(-x)=2^{-x}-2^{x}=-g(x)$  确为奇函数。

\textbf{方法二（直接用定义）}
\[f(-x)=(-x)^{3}\left(a\cdot2^{-x}-2^{x}\right)=-x^{3}\left(a\cdot2^{-x}-2^{x}\right)\]

由 $f(-x)=f(x)$：
\[-x^{3}\left(a\cdot2^{-x}-2^{x}\right)=x^{3}\left(a\cdot2^{x}-2^{-x}\right)\]

对 $x\ne0$ 约去 $x^{3}$：
\[-a\cdot2^{-x}+2^{x}=a\cdot2^{x}-2^{-x}\]
\[2^{x}+2^{-x}=a\left(2^{x}+2^{-x}\right)\]

因 $2^{x}+2^{-x}>0$，得 $a=1$。

\gkanswer{答案：$1$}

\begin{gktip}
\textbf{方法总结}：见到"$x^{\text{奇数次}}\times(\cdots)$ 是偶函数"，立刻用\textbf{奇偶性乘法规律}判定括号内是奇函数，再用 $g(0)=0$ 一步求出参数——比代定义展开快得多。同理，"偶 × 偶 = 偶"、"奇 × 偶 = 奇"也常用。
\end{gktip}

\par\medskip

\end{answer}
\begin{answer}
\gksolhead{14}{}

\textbf{题目}　抛物线 $C:y^{2}=2px\ (p>0)$ 的焦点为 $F$，$P$ 为 $C$ 上一点，$PF\perp x$ 轴，$Q$ 为 $x$ 轴上一点，且 $PQ\perp OP$。若 $|FQ|=6$，则 $C$ 的准线方程为 \gkblank{}。

\textbf{解析}

\textbf{第一步，定出 $P$ 的坐标。}

焦点 $F\left(\dfrac p2,0\right)$。由 $PF\perp x$ 轴知 $P$ 与 $F$ 横坐标相同：
\[x_P=\frac p2\]

代入抛物线方程：
\[y_P^{2}=2p\cdot\frac p2=p^{2}\ \Longrightarrow\ y_P=\pm p\]

不妨取 $P\left(\dfrac p2,\ p\right)$（取负号由对称性结果相同）。

\textbf{第二步，用垂直条件定出 $Q$。}

设 $Q(q,0)$。
\[\overrightarrow{OP}=\left(\frac p2,\ p\right),\qquad \overrightarrow{PQ}=\left(q-\frac p2,\ -p\right)\]

由 $PQ\perp OP$：
\[\overrightarrow{OP}\cdot\overrightarrow{PQ}=\frac p2\left(q-\frac p2\right)+p\times(-p)=0\]

\[\frac{pq}2-\frac{p^{2}}4-p^{2}=0\ \Longrightarrow\ \frac{pq}2=\frac{5p^{2}}4\]

因 $p>0$，两边除以 $\dfrac p2$：
\[q=\frac{5p}2\]

\textbf{第三步，由 $|FQ|=6$ 求 $p$。}
\[|FQ|=\left|q-\frac p2\right|=\left|\frac{5p}2-\frac p2\right|=2p=6\ \Longrightarrow\ p=3\]

\textbf{第四步，写出准线。}
\[x=-\frac p2=-\frac32\]

\gkanswer{答案：$x=-\dfrac32$}

\begin{gktip}
\textbf{验证}：$p=3$ 时 $P\left(\dfrac32,3\right)$，$Q\left(\dfrac{15}2,0\right)$，$\overrightarrow{OP}=\left(\dfrac32,3\right)$、$\overrightarrow{PQ}=(6,-3)$，数量积 $=9-9=0$  垂直成立；$|FQ|=\dfrac{15}2-\dfrac32=6$ 
\end{gktip}

\par\medskip

\end{answer}
\begin{answer}
\gksolhead{15}{}

\textbf{题目}　函数 $f(x)=|2x-1|-2\ln x$ 的最小值为 \gkblank{}。

\textbf{解析}

\textbf{定义域}：由 $\ln x$ 要求 $x>0$。

\textbf{按绝对值的零点 $x=\dfrac12$ 分段讨论。}

\par\medskip

\textbf{情形一：$x\ge\dfrac12$}
\[f(x)=2x-1-2\ln x\]
\[f'(x)=2-\frac2x=\frac{2(x-1)}x\]

\begin{itemize}
\item $\dfrac12\le x<1$：$f'<0$，递减
\item $x>1$：$f'>0$，递增
\end{itemize}

故在 $x=1$ 处取得该段最小值：
\[f(1)=2-1-2\ln1=1\]

\par\medskip

\textbf{情形二：$0<x<\dfrac12$}
\[f(x)=1-2x-2\ln x\]
\[f'(x)=-2-\frac2x<0\]

在 $\left(0,\dfrac12\right)$ 上\textbf{恒递减}，故该段的下确界在右端点处：
\[\lim_{x\to\frac12^{-}}f(x)=1-1-2\ln\frac12=2\ln2\approx1.386\]

该段的所有函数值都 $>2\ln2>1$。

\par\medskip

\textbf{综合两段}：全局最小值为
\[f(1)=1\]

\gkanswer{答案：$1$}

\begin{gktip}
\textbf{方法总结}：含绝对值的函数求最值，标准流程是"\textbf{按绝对值零点分段 → 每段分别求导找极值 → 比较各段最值与分界点的值}"。本题的关键是不要漏掉 $0<x<\dfrac12$ 这一段的分析（虽然它最终不产生最小值，但必须验证）。
\end{gktip}

\par\medskip

\end{answer}
\begin{answer}
\gksolhead{16}{}

\textbf{题目}　$20\ \mathrm{dm}\times12\ \mathrm{dm}$ 的长方形纸，沿对称轴对折。对折 $1$ 次得两种规格，面积和 $S_1=240$；对折 $2$ 次得三种规格，面积和 $S_2=180$。求对折 $4$ 次的规格种数与 $\displaystyle\sum_{k=1}^{n}S_k$。

\textbf{解析}

\textbf{第一步，找出规格的一般形式。}

每次对折可以沿长边或沿宽边。设对折 $n$ 次中，\textbf{长边被折了 $i$ 次}、宽边被折了 $n-i$ 次（$i=0,1,\cdots,n$），则规格为
\[\frac{20}{2^{i}}\ \mathrm{dm}\times\frac{12}{2^{\,n-i}}\ \mathrm{dm}\]

不同的 $i$ 给出不同的规格，故对折 $n$ 次共有
\[n+1\ \text{种规格}\]

\textbf{验证}：$n=1$ 时 $2$ 种（$10\times12$、$20\times6$）；$n=2$ 时 $3$ 种 

\textbf{对折 $4$ 次}：$4+1=5$ 种。

\textbf{第二步，求 $S_n$ 的表达式。}

\textbf{关键观察}：无论怎么折，对折 $n$ 次后每张纸的面积都是原面积的 $\dfrac1{2^{n}}$：
\[\text{每种规格的面积}=\frac{20\times12}{2^{n}}=\frac{240}{2^{n}}\]

共 $n+1$ 种规格，故
\[S_n=(n+1)\cdot\frac{240}{2^{n}}\]

\textbf{验证}：
\[S_1=2\times\frac{240}2=240\ \checkmark,\qquad S_2=3\times\frac{240}4=180\ \checkmark\]

\textbf{第三步，求和。}
\[\sum_{k=1}^{n}S_k=240\sum_{k=1}^{n}\frac{k+1}{2^{k}}=240\left(\underbrace{\sum_{k=1}^{n}\frac k{2^{k}}}_{A}+\underbrace{\sum_{k=1}^{n}\frac1{2^{k}}}_{B}\right)\]

\textbf{求 $B$（等比求和）}：
\[B=\frac{\frac12\left(1-\frac1{2^{n}}\right)}{1-\frac12}=1-\frac1{2^{n}}\]

\textbf{求 $A$（错位相减的经典结果）}：
\[A=\sum_{k=1}^{n}\frac k{2^{k}}=2-\frac{n+2}{2^{n}}\]

（推导：设 $A=\dfrac12+\dfrac2{4}+\cdots+\dfrac n{2^{n}}$，则 $\dfrac A2=\dfrac14+\dfrac2 8+\cdots+\dfrac n{2^{n+1}}$，相减得 $\dfrac A2=B-\dfrac n{2^{n+1}}$，解出即可。）

\textbf{合并}：
\[A+B=2-\frac{n+2}{2^{n}}+1-\frac1{2^{n}}=3-\frac{n+3}{2^{n}}\]

\[\sum_{k=1}^{n}S_k=240\left(3-\frac{n+3}{2^{n}}\right)=720-\frac{240(n+3)}{2^{n}}\]

\textbf{验证}：
\begin{itemize}
\item $n=1$：$720-\dfrac{240\times4}2=720-480=240=S_1$ 
\item $n=2$：$720-\dfrac{240\times5}4=720-300=420=240+180$ 
\end{itemize}

\gkanswer{答案：$5$ 种；$\displaystyle\sum_{k=1}^{n}S_k=720-\dfrac{240(n+3)}{2^{n}}$}

\begin{gktip}
\textbf{方法总结}：这是一道"\textbf{情境建模 + 错位相减}"的复合题。突破口有两个：① 意识到"折 $n$ 次的规格数 = 长边折的次数的可能取值个数 = $n+1$"；② 意识到"每张纸的面积恒为 $\dfrac{240}{2^{n}}$，与怎么折无关"。这两点想通后，$S_n=(n+1)\cdot\dfrac{240}{2^{n}}$ 立刻得出，剩下就是标准的 $\dfrac{\text{\text{一次式}}}{\text{\text{指数}}}$ 型求和。
\end{gktip}

\par\medskip

\end{answer}
\gkpart{四、解答题（共 70 分）}

\begin{answer}
\gksolhead{17}{10 分}

\textbf{题目}　$a_1=1$，$a_{n+1}=\begin{cases}a_n+1,&n\text{ \text{为奇数}}\\ a_n+2,&n\text{ \text{为偶数}}\end{cases}$

  (1) 记 $b_n=a_{2n}$，写出 $b_1$、$b_2$，并求 $\left\{b_n\right\}$ 的通项；(2) 求 $\left\{a_n\right\}$ 的前 $20$ 项和。

\par\medskip

\textbf{(1) 解}

\textbf{先逐项算出前几项}（注意条件中的 $n$ 指\textbf{下标}的奇偶）：
\[a_1=1\]
\[a_2=a_1+1=2\quad(n=1\ \text{奇})\]
\[a_3=a_2+2=4\quad(n=2\ \text{偶})\]
\[a_4=a_3+1=5\quad(n=3\ \text{奇})\]
\[a_5=a_4+2=7,\qquad a_6=a_5+1=8\]

故
\[b_1=a_2=2,\qquad b_2=a_4=5\]

\textbf{求通项}：把"两步"合并成一步。
\[a_{2n+1}=a_{2n}+2\quad(\text{下标 }2n\text{ 为偶})\]
\[a_{2n+2}=a_{2n+1}+1\quad(\text{下标 }2n+1\text{ 为奇})\]

两式相加消去中间项：
\[a_{2n+2}=a_{2n}+3\ \Longrightarrow\ b_{n+1}=b_n+3\]

$\left\{b_n\right\}$ 是首项 $b_1=2$、公差 $3$ 的\textbf{等差数列}：
\[b_n=2+3(n-1)=3n-1\]

\textbf{检验}：$b_2=5$ ，$b_3=8=a_6$ 

\par\medskip

\textbf{(2) 解}

前 $20$ 项分成奇数项与偶数项两组分别求和。

\textbf{偶数项}（$a_2,a_4,\cdots,a_{20}$，即 $b_1\sim b_{10}$）：
\[\sum_{n=1}^{10}b_n=\sum_{n=1}^{10}(3n-1)=3\times\frac{10\times11}2-10=165-10=155\]

\textbf{奇数项}（$a_1,a_3,\cdots,a_{19}$，共 $10$ 项）：

由 $a_{2n+1}=a_{2n}+2=b_n+2=3n+1$：
\begin{itemize}
\item $n=1$：$a_3=4$
\item $n=2$：$a_5=7$
\item $\cdots$
\item $n=9$：$a_{19}=28$
\end{itemize}

再加上单独的 $a_1=1$：
\[a_1+\sum_{n=1}^{9}(3n+1)=1+\left(3\times\frac{9\times10}2+9\right)=1+(135+9)=145\]

\textbf{总和}：
\[S_{20}=145+155=300\]

\gkanswer{答案：$300$}

\begin{gktip}
\textbf{方法总结}：分奇偶定义的递推数列，通法是"\textbf{隔项合并}"——把相邻两步的递推相加，消去中间项，得到 $a_{n+2}$ 与 $a_n$ 的关系，从而把奇数项、偶数项各自变成简单数列（本题都成了等差数列）。求和时也按奇偶分组，注意\textbf{首项 $a_1$ 是否要单独处理}（本题奇数项的通项公式从 $a_3$ 开始才适用）。
\end{gktip}

\par\medskip

\end{answer}
\begin{answer}
\gksolhead{18}{12 分}

\textbf{题目}　$A$ 类问题答对得 $20$ 分、$B$ 类答对得 $80$ 分，答错得 $0$ 分并结束；答对则从另一类再抽一题，之后无论对错都结束。小明答对 $A$ 的概率 $0.8$、答对 $B$ 的概率 $0.6$。

  (1) 若先答 $A$，求累计得分 $X$ 的分布列；(2) 为使期望最大，应先答哪类？

\par\medskip

\textbf{(1) 解}

\textbf{先理清流程}（先答 $A$）：

\begin{itemize}
\item 第一题（$A$ 类）答\textbf{错}：得 $0$ 分，比赛结束
\item 第一题答\textbf{对}：得 $20$ 分，继续答 $B$ 类
\item $B$ 答\textbf{对}：再得 $80$ 分，总分 $100$
\item $B$ 答\textbf{错}：总分仍为 $20$
\end{itemize}

\textbf{故 $X$ 的可能取值为 $0$、$20$、$100$}，各自概率：

\[P(X=0)=1-0.8=0.2\]
\[P(X=20)=0.8\times(1-0.6)=0.8\times0.4=0.32\]
\[P(X=100)=0.8\times0.6=0.48\]

\textbf{分布列}

\begin{center}
\begin{tabular}{cccc}
\toprule
$X$ & $0$ & $20$ & $100$ \\
\midrule
$P$ & $0.2$ & $0.32$ & $0.48$ \\
\bottomrule
\end{tabular}
\end{center}

\textbf{校验}：$0.2+0.32+0.48=1$ 

\textbf{期望}：
\[E(X)=0\times0.2+20\times0.32+100\times0.48=6.4+48=54.4\]

\par\medskip

\textbf{(2) 解}

\textbf{计算先答 $B$ 的期望。} 设累计得分为 $Y$：

\begin{itemize}
\item $B$ 答错：$Y=0$，概率 $1-0.6=0.4$
\item $B$ 答对（得 $80$ 分）后答 $A$：
\item $A$ 答错：$Y=80$，概率 $0.6\times0.2=0.12$
\item $A$ 答对：$Y=100$，概率 $0.6\times0.8=0.48$
\end{itemize}

\begin{center}
\begin{tabular}{cccc}
\toprule
$Y$ & $0$ & $80$ & $100$ \\
\midrule
$P$ & $0.4$ & $0.12$ & $0.48$ \\
\bottomrule
\end{tabular}
\end{center}

\[E(Y)=0\times0.4+80\times0.12+100\times0.48=9.6+48=57.6\]

\textbf{比较}：
\[E(Y)=57.6>54.4=E(X)\]

\textbf{所以小明应选择先回答 $B$ 类问题。}

\par\medskip

\textbf{理由的直观解释}（答题时写上会更完整）：

两种顺序下，"两题全对得 $100$ 分"的概率\textbf{完全相同}（都是 $0.8\times0.6=0.48$，乘法交换律）。差别只在于"\textbf{第一题对、第二题错}"这种情形能保住多少分：

\begin{itemize}
\item 先答 $B$：保住 $80$ 分，概率 $0.12$，贡献 $9.6$
\item 先答 $A$：保住 $20$ 分，概率 $0.32$，贡献 $6.4$
\end{itemize}

虽然先答 $A$ 时"保住分数"的概率更大（$0.32>0.12$），但保住的分值太低（$20$ 对 $80$），总体上不划算。

\textbf{一般原则}：这类"闯关式"问题，\textbf{应该把分值高的题放在前面}（只要其正确率不是低得离谱）。

\begin{gktip}
\textbf{方法总结}：分布列题的关键是\textbf{理清流程图}——先画出决策树，标清每条路径的概率与得分，再合并同类得分。本题第 (2) 问只需算另一种顺序的期望并比较，但\textbf{必须写出理由}（很多学生只算不说，会被扣分）。
\end{gktip}

\par\medskip

\end{answer}
\begin{answer}
\gksolhead{19}{12 分}

\textbf{题目}　$\triangle ABC$ 中 $b^{2}=ac$，$D$ 在边 $AC$ 上，$BD\sin\angle ABC=a\sin C$。(1) 证明 $BD=b$；(2) 若 $AD=2DC$，求 $\cos\angle ABC$。

\par\medskip

\textbf{(1) 证明}

记 $\angle ABC=B$。由\textbf{正弦定理}（设外接圆半径为 $R$）：
\[\sin B=\frac b{2R},\qquad \sin C=\frac c{2R}\]

代入已知条件 $BD\sin B=a\sin C$：
\[BD\cdot\frac b{2R}=a\cdot\frac c{2R}\]

两边同乘 $2R$：
\[BD\cdot b=ac\]

由题设 $b^{2}=ac$：
\[BD\cdot b=b^{2}\]

因 $b>0$，两边约去 $b$：
\[BD=b\qquad\blacksquare\]

\par\medskip

\textbf{(2) 解}

\textbf{第一步，标出各线段长度。}

$D$ 在 $AC$ 上且 $AD=2DC$，而 $AC=b$：
\[AD=\frac{2b}3,\qquad DC=\frac b3\]

由 (1)：$BD=b$。

\textbf{第二步，在两个小三角形中分别用余弦定理。}

设 $\angle BDA=\theta$，则 $\angle BDC=\pi-\theta$（\textbf{互补}，这是消元的关键）。

\textbf{在 $\triangle ABD$ 中}（边 $AD$、$BD$ 夹角 $\theta$，对边 $AB=c$）：
\[c^{2}=AD^{2}+BD^{2}-2\cdot AD\cdot BD\cos\theta=\frac{4b^{2}}9+b^{2}-\frac{4b^{2}}3\cos\theta\]

\[c^{2}=\frac{13b^{2}}9-\frac{4b^{2}}3\cos\theta\qquad(\text{i})\]

\textbf{在 $\triangle CBD$ 中}（边 $DC$、$BD$ 夹角 $\pi-\theta$，对边 $BC=a$，注意 $\cos(\pi-\theta)=-\cos\theta$）：
\[a^{2}=\frac{b^{2}}9+b^{2}+\frac{2b^{2}}3\cos\theta=\frac{10b^{2}}9+\frac{2b^{2}}3\cos\theta\qquad(\text{ii})\]

\textbf{第三步，消去 $\cos\theta$。}

观察 (i) 中 $\cos\theta$ 的系数为 $-\dfrac{4b^{2}}3$、(ii) 中为 $+\dfrac{2b^{2}}3$，按 $(\text{i})+2\times(\text{ii})$ 恰好抵消：
\[c^{2}+2a^{2}=\frac{13b^{2}}9+\frac{20b^{2}}9=\frac{33b^{2}}9=\frac{11b^{2}}3\]

再用 $b^{2}=ac$：
\[c^{2}+2a^{2}=\frac{11ac}3\]

\textbf{第四步，化为关于 $\dfrac ca$ 的方程。}

两边除以 $ac$：
\[\frac ca+\frac{2a}c=\frac{11}3\]

令 $t=\dfrac ca>0$：
\[t+\frac2t=\frac{11}3\ \Longrightarrow\ 3t^{2}-11t+6=0\]

\[t=\frac{11\pm\sqrt{121-72}}6=\frac{11\pm7}6\]

\[t=3\qquad\text{或}\qquad t=\frac23\]

\textbf{第五步，求 $\cos B$ 并筛选。}

由余弦定理及 $b^{2}=ac$：
\[\cos B=\frac{a^{2}+c^{2}-b^{2}}{2ac}=\frac{a^{2}+c^{2}-ac}{2ac}=\frac{1+t^{2}-t}{2t}\]

\begin{itemize}
\item $t=3$：$\cos B=\dfrac{1+9-3}6=\dfrac76>1$  \textbf{不可能}（余弦值不能超过 $1$），舍去
\item $t=\dfrac23$：
\end{itemize}
\[\cos B=\frac{1+\dfrac49-\dfrac23}{\dfrac43}=\frac{\dfrac{9+4-6}9}{\dfrac43}=\frac{\dfrac79}{\dfrac43}=\frac79\times\frac34=\frac7{12}\]

\textbf{所以 $\cos\angle ABC=\dfrac7{12}$。}

\begin{gktip}
\textbf{方法总结}：出现"三角形内的一条线段（如中线、角平分线、任意分点）"时，标准做法是\textbf{在两个小三角形中对互补角分别用余弦定理，再加权相加消去 $\cos\theta$}。权重的选取原则：让 $\cos\theta$ 的系数相消（本题是 $1:2$）。

最后\textbf{必须筛选}：$t=3$ 给出的 $\cos B>1$ 在数学上不可能，这类"增根"是命题人设置的检验点。
\end{gktip}

\par\medskip

\end{answer}
\begin{answer}
\gksolhead{20}{12 分}

\textbf{题目}　三棱锥 $A\text{-}BCD$ 中，平面 $ABD\perp$ 平面 $BCD$，$AB=AD$，$O$ 为 $BD$ 中点。

\begin{center}% 2021 新高考Ⅰ卷 20：三棱锥 A-BCD（平面 ABD⊥平面 BCD，O 为 BD 中点，E 在 AD 上）
\begin{tikzpicture}[scale=1.0, font=\small]
  \coordinate (A) at (2.1,2.4);
  \coordinate (B) at (0,0.75);
  \coordinate (C) at (2.35,0);
  \coordinate (D) at (4.35,0.8);
  \coordinate (O) at ($(B)!0.5!(D)$);
  \coordinate (E) at ($(A)!0.3333!(D)$);
  \draw (A) -- (B) -- (C) -- cycle;
  \draw (A) -- (D)  (C) -- (D);
  \draw[dashed] (B) -- (D);
  \draw[dashed] (A) -- (O)  (C) -- (O);
  \draw (C) -- (E);
  \fill (O) circle (0.9pt);
  \node[above] at (A) {$A$};
  \node[left] at (B) {$B$};
  \node[below] at (C) {$C$};
  \node[right] at (D) {$D$};
  \node[above left, inner sep=0.5pt] at (O) {$O$};
  \node[above right, inner sep=1pt] at (E) {$E$};
\end{tikzpicture}
\end{center}

  (1) 证明 $OA\perp CD$；(2) 若 $\triangle OCD$ 是边长为 $1$ 的等边三角形，$E$ 在棱 $AD$ 上且 $DE=2EA$，二面角 $E\text{-}BC\text{-}D$ 为 $45^\circ$，求三棱锥 $A\text{-}BCD$ 的体积。

\par\medskip

\textbf{(1) 证明}

因为 $AB=AD$ 且 $O$ 是 $BD$ 的中点，由\textbf{等腰三角形三线合一}：
\[AO\perp BD\]

又平面 $ABD\perp$ 平面 $BCD$，两平面的\textbf{交线为 $BD$}，而 $AO\subset$ 平面 $ABD$ 且 $AO\perp BD$，

由\textbf{面面垂直的性质定理}：
\[AO\perp\text{平面 }BCD\]

又 $CD\subset$ 平面 $BCD$，由线面垂直的定义：
\[OA\perp CD\qquad\blacksquare\]

\par\medskip

\textbf{(2) 解}

\textbf{建系}：由 (1) 知 $AO\perp$ 平面 $BCD$，故以 $O$ 为原点、$OA$ 为 $z$ 轴建立空间直角坐标系。

由 $\triangle OCD$ 是边长 $1$ 的等边三角形，$OD=1$；又 $O$ 是 $BD$ 中点，故 $OB=1$、$BD=2$。

把 $D$ 放在 $x$ 轴正方向：
\[O(0,0,0),\quad D(1,0,0),\quad B(-1,0,0)\]

$C$ 满足 $OC=CD=1$，即 $\triangle OCD$ 等边，故
\[C\left(\frac12,\ \frac{\sqrt3}2,\ 0\right)\]

设 $A(0,0,h)$（$h=OA>0$，待求）。

\textbf{求 $E$ 的坐标}：$E$ 在棱 $AD$ 上且 $DE=2EA$，故 $E$ 分 $DA$ 为 $2:1$：
\[E=D+\frac23\left(A-D\right)=(1,0,0)+\frac23(-1,0,h)=\left(\frac13,\ 0,\ \frac{2h}3\right)\]

\textbf{求两个半平面的法向量}（棱为 $BC$）：

\[\overrightarrow{BC}=C-B=\left(\frac32,\ \frac{\sqrt3}2,\ 0\right),\qquad \overrightarrow{BE}=E-B=\left(\frac43,\ 0,\ \frac{2h}3\right)\]

\textbf{平面 $EBC$} 的法向量：
\[\boldsymbol n_1=\overrightarrow{BC}\times\overrightarrow{BE}=\left(\frac{\sqrt3}2\cdot\frac{2h}3-0,\ \ 0-\frac32\cdot\frac{2h}3,\ \ 0-\frac{\sqrt3}2\cdot\frac43\right)=\left(\frac{\sqrt3h}3,\ -h,\ -\frac{2\sqrt3}3\right)\]

\textbf{平面 $BCD$} 就是 $z=0$ 平面，法向量：
\[\boldsymbol n_2=(0,0,1)\]

\textbf{由二面角为 $45^\circ$ 列方程}：
\[\left|\boldsymbol n_1\right|^{2}=\frac{3h^{2}}9+h^{2}+\frac{12}9=\frac{h^{2}}3+h^{2}+\frac43=\frac{4\left(h^{2}+1\right)}3\]

\[\cos45^\circ=\frac{\left|\boldsymbol n_1\cdot\boldsymbol n_2\right|}{\left|\boldsymbol n_1\right|\left|\boldsymbol n_2\right|}=\frac{\dfrac{2\sqrt3}3}{\dfrac2{\sqrt3}\sqrt{h^{2}+1}}=\frac{\dfrac{2\sqrt3}3\cdot\sqrt3}{2\sqrt{h^{2}+1}}=\frac{1}{\sqrt{h^{2}+1}}\]

令其等于 $\dfrac{\sqrt2}2$：
\[\frac1{\sqrt{h^{2}+1}}=\frac1{\sqrt2}\ \Longrightarrow\ h^{2}+1=2\ \Longrightarrow\ h=1\]

\textbf{求体积}：

底面 $\triangle BCD$：底边 $BD=2$（在 $x$ 轴上），$C$ 到 $x$ 轴的距离为 $\dfrac{\sqrt3}2$：
\[S_{\triangle BCD}=\frac12\times2\times\frac{\sqrt3}2=\frac{\sqrt3}2\]

高为 $OA=h=1$：
\[V=\frac13\times\frac{\sqrt3}2\times1=\frac{\sqrt3}6\]

\textbf{所以三棱锥 $A\text{-}BCD$ 的体积为 $\dfrac{\sqrt3}6$。}

\begin{gktip}
\textbf{方法总结}：本题的建系依据完全来自第 (1) 问的结论——$AO\perp$ 平面 $BCD$ 使得 $OA$ 可以直接作为 $z$ 轴。\textbf{在解答题中，第 (1) 问往往是第 (2) 问建系的"许可证"}，务必先做 (1) 再做 (2)。

另外，二面角的一个面若是\textbf{坐标平面}（本题的 $z=0$），其法向量直接取 $(0,0,1)$，可省去一次叉乘。
\end{gktip}

\par\medskip

\end{answer}
\begin{answer}
\gksolhead{21}{12 分}

\textbf{题目}　$F_1\left(-\sqrt{17},0\right)$、$F_2\left(\sqrt{17},0\right)$，$M$ 满足 $\left|MF_1\right|-\left|MF_2\right|=2$，轨迹为 $C$。(1) 求 $C$ 的方程；(2) $T$ 在直线 $x=\dfrac12$ 上，过 $T$ 的两条直线分别交 $C$ 于 $A,B$ 与 $P,Q$，且 $|TA|\cdot|TB|=|TP|\cdot|TQ|$，求 $k_{AB}+k_{PQ}$。

\par\medskip

\textbf{(1) 解}

\textbf{识别曲线类型}：到两定点的距离\textbf{之差为定值}（且该定值小于两定点距离 $2\sqrt{17}$），故轨迹是双曲线。

注意差是 $\left|MF_1\right|-\left|MF_2\right|=2>0$（\textbf{没有加绝对值}），说明 $M$ 到左焦点更远，即轨迹只是\textbf{右支}。

由定义：
\[2a=2\ \Longrightarrow\ a=1\]
\[c=\sqrt{17}\ \Longrightarrow\ b^{2}=c^{2}-a^{2}=17-1=16\]

\[C:\ x^{2}-\frac{y^{2}}{16}=1\qquad(x\ge1)\]

\par\medskip

\textbf{(2) 解}

\textbf{第一步，用直线的参数方程统一处理（本题最简洁的路线）。}

设 $T\left(\dfrac12,\ t\right)$。过 $T$ 且倾斜角为 $\theta$ 的直线的参数方程为
\[\begin{cases}x=\dfrac12+s\cos\theta\\ y=t+s\sin\theta\end{cases}\]

其中参数 $s$ 表示点到 $T$ 的\textbf{有向距离}——这正是本题需要的量。

\textbf{第二步，代入曲线方程。}

把双曲线写成 $16x^{2}-y^{2}=16$：
\[16\left(\frac12+s\cos\theta\right)^{2}-\left(t+s\sin\theta\right)^{2}=16\]

展开：
\[16\left(\frac14+s\cos\theta+s^{2}\cos^{2}\theta\right)-\left(t^{2}+2ts\sin\theta+s^{2}\sin^{2}\theta\right)=16\]

整理成关于 $s$ 的二次方程：
\[\left(16\cos^{2}\theta-\sin^{2}\theta\right)s^{2}+\left(16\cos\theta-2t\sin\theta\right)s-\left(12+t^{2}\right)=0\]

\textbf{第三步，由韦达定理写出距离之积。}
\[|TA|\cdot|TB|=\left|s_1s_2\right|=\left|\frac{-\left(12+t^{2}\right)}{16\cos^{2}\theta-\sin^{2}\theta}\right|=\frac{12+t^{2}}{\left|16\cos^{2}\theta-\sin^{2}\theta\right|}\]

把 $\cos^2\theta$ 换成斜率 $k=\tan\theta$（用 $\cos^{2}\theta=\dfrac1{1+k^{2}}$）：
\[16\cos^{2}\theta-\sin^{2}\theta=\cos^{2}\theta\left(16-k^{2}\right)=\frac{16-k^{2}}{1+k^{2}}\]

\[|TA|\cdot|TB|=\frac{\left(12+t^{2}\right)\left(1+k^{2}\right)}{\left|16-k^{2}\right|}\]

\textbf{第四步，代入条件。}

记 $k_1=k_{AB}$、$k_2=k_{PQ}$。由 $|TA||TB|=|TP||TQ|$（注意 $12+t^{2}$ 是两边公共的因子，可约去）：
\[\frac{1+k_1^{2}}{\left|16-k_1^{2}\right|}=\frac{1+k_2^{2}}{\left|16-k_2^{2}\right|}\]

\textbf{第五步（关键）：判断绝对值内的符号。}

双曲线 $x^{2}-\dfrac{y^{2}}{16}=1$ 的渐近线斜率为 $\pm4$。

点 $T\left(\dfrac12,t\right)$ 满足 $\dfrac14-\dfrac{t^{2}}{16}<1$，说明 $T$ 位于双曲线的\textbf{中心区域}（两支之间）。从这样的点出发的直线要与\textbf{同一支}相交于两点，其斜率的绝对值必须\textbf{大于渐近线斜率}：
\[\left|k_1\right|>4,\qquad \left|k_2\right|>4\]

故 $16-k_1^{2}<0$ 且 $16-k_2^{2}<0$，两个绝对值内\textbf{同号}，可以同时去掉绝对值（负负相消）：
\[\frac{1+k_1^{2}}{16-k_1^{2}}=\frac{1+k_2^{2}}{16-k_2^{2}}\]

\textbf{第六步，化简。}

交叉相乘：
\[\left(1+k_1^{2}\right)\left(16-k_2^{2}\right)=\left(1+k_2^{2}\right)\left(16-k_1^{2}\right)\]

\[16-k_2^{2}+16k_1^{2}-k_1^{2}k_2^{2}=16-k_1^{2}+16k_2^{2}-k_1^{2}k_2^{2}\]

消去公共项 $16$ 与 $-k_1^{2}k_2^{2}$：
\[-k_2^{2}+16k_1^{2}=-k_1^{2}+16k_2^{2}\ \Longrightarrow\ 17k_1^{2}=17k_2^{2}\ \Longrightarrow\ k_1^{2}=k_2^{2}\]

因两条直线不同（$k_1\ne k_2$），必有
\[k_1=-k_2\]

\[k_{AB}+k_{PQ}=0\]

\gkanswer{答案：斜率之和为 $0$}

\begin{gktip}
\textbf{方法总结}：本题用\textbf{直线参数方程}是最优解——因为条件里的 $|TA|\cdot|TB|$ 正是参数 $s$ 的两根之积的绝对值，用普通方程需要先求交点坐标再算距离，麻烦得多。

第五步的符号判断是本题最容易出错的地方：若不论证"$|k|>4$"，就必须讨论绝对值内异号的情形，会引入多余的解。\textbf{从双曲线内部的点出发的直线要与同一支交两点，斜率必须比渐近线更陡}——这个几何事实值得记住。
\end{gktip}

\par\medskip

\end{answer}
\begin{answer}
\gksolhead{22}{12 分}

\textbf{题目}　已知 $f(x)=x(1-\ln x)$。(1) 讨论单调性；(2) 设 $a$、$b$ 为不相等的正数且 $b\ln a-a\ln b=a-b$，证明 $2<\dfrac1a+\dfrac1b<\mathrm e$。

\par\medskip

\textbf{(1) 解}

定义域：$x>0$。

\textbf{求导}（乘积法则）：
\[f'(x)=1\cdot(1-\ln x)+x\cdot\left(-\frac1x\right)=1-\ln x-1=-\ln x\]

\textbf{这个导数出奇地简洁}，是本题设计的巧妙之处。

\begin{itemize}
\item $0<x<1$ 时 $\ln x<0$，故 $f'(x)>0$，$f$ \textbf{单调递增}
\item $x>1$ 时 $\ln x>0$，故 $f'(x)<0$，$f$ \textbf{单调递减}
\end{itemize}

\textbf{最大值}：
\[f(1)=1\times(1-0)=1\]

\textbf{其它有用的信息}（第 (2) 问要用）：
\begin{itemize}
\item $f(\mathrm e)=\mathrm e(1-1)=0$
\item $f''(x)=-\dfrac1x<0$，故 $f$ 是\textbf{凹函数}（图象上凸）
\end{itemize}

\par\medskip

\textbf{(2) 证明}

\textbf{第一步（本题的钥匙）：换元，把条件化成 $f(u)=f(v)$。}

条件为
\[b\ln a-a\ln b=a-b\]

两边同除以 $ab$（$a,b>0$）：
\[\frac{\ln a}a-\frac{\ln b}b=\frac1b-\frac1a\]

移项，把下标相同的归到一边：
\[\frac{\ln a}a+\frac1a=\frac{\ln b}b+\frac1b\ \Longrightarrow\ \frac{1+\ln a}a=\frac{1+\ln b}b\]

\textbf{令}
\[u=\frac1a,\qquad v=\frac1b\]

则 $a=\dfrac1u$，$\ln a=-\ln u$，于是
\[\frac{1+\ln a}a=u\left(1-\ln u\right)=f(u)\]

同理右边 $=f(v)$。\textbf{条件化为}
\[\boxed{f(u)=f(v),\qquad u\ne v,\ u,v>0}\]

\textbf{而所要证明的正是} $2<u+v<\mathrm e$。

由 (1) 的单调性（$f$ 在 $(0,1)$ 增、$(1,+\infty)$ 减），$f(u)=f(v)$ 且 $u\ne v$ 意味着两者分居 $1$ 的两侧，不妨设
\[0<u<1<v\]

\par\medskip

\textbf{第二步，证明左半 $u+v>2$。}

\textbf{思路}：证 $v>2-u$。因 $2-u\in(1,2)$ 与 $v$ 都在 $f$ 的\textbf{递减}区间上，只需证
\[f(v)<f(2-u)\]

而 $f(v)=f(u)$，故只需证
\[f(u)<f(2-u)\qquad\text{对 }u\in(0,1)\]

\textbf{构造函数}：
\[g(x)=f(x)-f(2-x),\qquad x\in(0,1)\]

\[g'(x)=f'(x)+f'(2-x)=-\ln x-\ln(2-x)=-\ln\left[x(2-x)\right]\]

因
\[x(2-x)=1-(x-1)^{2}<1\qquad(x\ne1)\]

故 $\ln\left[x(2-x)\right]<0$，从而
\[g'(x)>0\]

$g$ 在 $(0,1)$ 上\textbf{严格递增}，且
\[g(1)=f(1)-f(1)=0\]

所以对 $x\in(0,1)$ 有 $g(x)<0$，即
\[f(u)<f(2-u)\quad\Longrightarrow\quad f(v)<f(2-u)\]

由 $f$ 在 $(1,+\infty)$ 上递减，得
\[v>2-u\ \Longrightarrow\ u+v>2\qquad\checkmark\]

\par\medskip

\textbf{第三步，证明右半 $u+v<\mathrm e$。}

\textbf{思路：用凹函数在 $x=\mathrm e$ 处的切线放缩。}

由 (1)，$f''(x)=-\dfrac1x<0$，$f$ 是凹函数，故\textbf{图象位于任意一条切线的下方}。

在 $x=\mathrm e$ 处：$f(\mathrm e)=0$，$f'(\mathrm e)=-\ln\mathrm e=-1$，切线为
\[y=0-1\cdot(x-\mathrm e)=\mathrm e-x\]

于是
\[f(x)\le\mathrm e-x\qquad\text{对一切 }x>0,\ \text{等号仅在 }x=\mathrm e\]

\textbf{应用于 $v$}：因 $f(v)=f(u)>0$（$u\in(0,1)$ 时 $1-\ln u>0$，故 $f(u)>0$），而 $f(\mathrm e)=0$，故 $v\ne\mathrm e$，切线不等式\textbf{严格成立}：
\[f(v)<\mathrm e-v\ \Longrightarrow\ v<\mathrm e-f(v)=\mathrm e-f(u)\]

\textbf{代入 $f(u)=u(1-\ln u)$}：
\[u+v<u+\mathrm e-u(1-\ln u)=u+\mathrm e-u+u\ln u=\mathrm e+u\ln u\]

\textbf{最后一步}：因 $0<u<1$，有 $\ln u<0$，故 $u\ln u<0$：
\[u+v<\mathrm e+u\ln u<\mathrm e\qquad\checkmark\]

\par\medskip

\textbf{综上}：
\[2<\frac1a+\frac1b=u+v<\mathrm e\qquad\blacksquare\]

\begin{gktip}
\textbf{方法总结}：这道压轴题的三个关键点值得反复体会——
1. \textbf{换元识别结构}：条件 $b\ln a-a\ln b=a-b$ 看似陌生，除以 $ab$ 并令 $u=\dfrac1a$、$v=\dfrac1b$ 后，竟然恰好化为第 (1) 问那个函数的 $f(u)=f(v)$。\textbf{第 (1) 问从来不是孤立的，它是第 (2) 问的工具}。
2. \textbf{左半用"对称化构造"}：证 $u+v>2$（即两根之和大于两倍极值点）是\textbf{极值点偏移}问题的标准形式，通法就是构造 $g(x)=f(x)-f(2x_0-x)$ 并判断其单调性。
3. \textbf{右半用"切线放缩"}：凹函数在任意点的切线都在图象上方，选取合适的切点（本题选 $x=\mathrm e$ 因为 $f(\mathrm e)=0$ 最简洁）可以把超越不等式转化为代数不等式。

\textbf{考试策略}：第 (1) 问必得分；第 (2) 问的换元若想不到，后面全部无法进行——这正是压轴题的区分点所在。
\end{gktip}


\end{answer}